
\documentclass{article}
\usepackage{colortbl}
\usepackage{makecell}
\usepackage{multirow}
\usepackage{supertabular}

\begin{document}

\newcounter{utterance}

\centering \large Interaction Transcript for game `matchit\_ascii', experiment `different\_grid', episode 3 with nemotron{-}3{-}ultra.
\vspace{24pt}

{ \footnotesize  \setcounter{utterance}{1}
\setlength{\tabcolsep}{0pt}
\begin{supertabular}{c@{$\;$}|p{.15\linewidth}@{}p{.15\linewidth}p{.15\linewidth}p{.15\linewidth}p{.15\linewidth}p{.15\linewidth}}
    \# & $\;$A & \multicolumn{4}{c}{Game Master} & $\;\:$B\\
    \hline

    \theutterance \stepcounter{utterance}  
    & & \multicolumn{4}{p{0.6\linewidth}}{
        \cellcolor[rgb]{0.9,0.9,0.9}{
            \makecell[{{p{\linewidth}}}]{
                \texttt{\tiny{[P1$\langle$GM]}}
                \texttt{Du nimmst an einem gemeinsamen Ratespiel teil.} \\
\texttt{Das Ziel ist herauszufinden, ob dieses Raster und ein anderes Bild, das nur ich sehen kann, gleich sind.} \\
\texttt{Bitte beschreibe zuerst dein Raster. Dann gebe ich meine Beschreibung ab und wir können uns gegenseitig Fragen zu den Rastern stellen, um herauszufinden, ob sie gleich sind. Beginne jetzt deine kurze Rasterbeschreibung mit "BESCHREIBUNG:" gefolgt von der Beschreibung. Füge nichts anderes hinzu. Das ist dein Raster:} \\
\texttt{▢ X X X ▢} \\
\texttt{X ▢ X ▢ X} \\
\texttt{X X ▢ X X} \\
\texttt{X X ▢ X X} \\
\texttt{X X ▢ X X} \\
            }
        }
    }
    & \\ \\

    \theutterance \stepcounter{utterance}  
    & \multicolumn{4}{p{0.6\linewidth}}{
        \cellcolor[rgb]{0.9,0.9,0.9}{
            \makecell[{{p{\linewidth}}}]{
                \texttt{\tiny{[P1$\rangle$GM]}}
                \texttt{BESCHREIBUNG: 5x5{-}Raster. Kreuzmuster in der Mitte: Zeile 3, Spalte 3 ist leer (▢). Zeile 1: Spalten 1, 4 leer. Zeile 2: Spalten 2, 4 leer. Zeilen 3{-}5: Spalte 3 leer. Übrige Felder sind X.} \\
            }
        }
    }
    & & \\ \\

    \theutterance \stepcounter{utterance}  
    & & & \multicolumn{2}{p{0.3\linewidth}}{
        \cellcolor[rgb]{0.9,0.9,0.9}{
            \makecell[{{p{\linewidth}}}]{
                \texttt{\tiny{[GM$|$GM]}}
                \texttt{continue} \\
            }
        }
    }
    & & \\ \\

    \theutterance \stepcounter{utterance}  
    & & \multicolumn{4}{p{0.6\linewidth}}{
        \cellcolor[rgb]{0.9,0.9,0.9}{
            \makecell[{{p{\linewidth}}}]{
                \texttt{\tiny{[GM$\rangle$P2]}}
                \texttt{Du nimmst an einem gemeinsamen Ratespiel teil.} \\
\texttt{Das Ziel ist herauszufinden, ob dieses Raster und ein anderes Bild, das nur ich sehen kann, gleich sind.} \\
\texttt{Bitte beschreibe zuerst dein Raster. Dann gebe ich meine Beschreibung ab und wir können uns gegenseitig Fragen zu den Rastern stellen, um herauszufinden, ob sie gleich sind. Beginne jetzt deine kurze Rasterbeschreibung mit "BESCHREIBUNG:" gefolgt von der Beschreibung. Füge nichts anderes hinzu. Das ist dein Raster:} \\
\texttt{▢ ▢ ▢ ▢ X} \\
\texttt{▢ ▢ ▢ X X} \\
\texttt{▢ ▢ X X X} \\
\texttt{▢ X X X X} \\
\texttt{X X X X X} \\
            }
        }
    }
    & \\ \\

    \theutterance \stepcounter{utterance}  
    & & & \multicolumn{4}{p{0.6\linewidth}}{
        \cellcolor[rgb]{0.9,0.9,0.9}{
            \makecell[{{p{\linewidth}}}]{
                \texttt{\tiny{[GM$\langle$P2]}}
                \texttt{BESCHREIBUNG: 5x5{-}Gitter. Diagonale von oben rechts nach unten links und alles darunter (bzw. links unten) ist mit X gefüllt. Zeile 1: 4 leer, 1 X (ganz rechts). Zeile 2: 3 leer, 2 X. Zeile 3: 2 leer, 3 X. Zeile 4: 1 leer, 4 X. Zeile 5: 5 X.} \\
            }
        }
    }
     \\ \\

    \theutterance \stepcounter{utterance}  
    & & & \multicolumn{2}{p{0.3\linewidth}}{
        \cellcolor[rgb]{0.9,0.9,0.9}{
            \makecell[{{p{\linewidth}}}]{
                \texttt{\tiny{[GM$|$GM]}}
                \texttt{continue} \\
            }
        }
    }
    & & \\ \\

    \theutterance \stepcounter{utterance}  
    & & \multicolumn{4}{p{0.6\linewidth}}{
        \cellcolor[rgb]{0.9,0.9,0.9}{
            \makecell[{{p{\linewidth}}}]{
                \texttt{\tiny{[GM$\rangle$P2]}}
                \texttt{Das ist meineBESCHREIBUNG: 5x5{-}Raster. Kreuzmuster in der Mitte: Zeile 3, Spalte 3 ist leer (▢). Zeile 1: Spalten 1, 4 leer. Zeile 2: Spalten 2, 4 leer. Zeilen 3{-}5: Spalte 3 leer. Übrige Felder sind X.} \\
\texttt{Stelle jetzt eine Frage, um neue Aspekte meines Rasters herauszufinden, die sich möglicherweise von deinem Raster unterscheiden. Beginne mit "FRAGE:" und füge nichts anderes hinzu.} \\
            }
        }
    }
    & \\ \\

    \theutterance \stepcounter{utterance}  
    & & & \multicolumn{4}{p{0.6\linewidth}}{
        \cellcolor[rgb]{0.9,0.9,0.9}{
            \makecell[{{p{\linewidth}}}]{
                \texttt{\tiny{[GM$\langle$P2]}}
                \texttt{FRAGE: Welchen Inhalt haben die Ecken des Rasters (Positionen Zeile 1 Spalte 1, Zeile 1 Spalte 5, Zeile 5 Spalte 1, Zeile 5 Spalte 5)?} \\
            }
        }
    }
     \\ \\

    \theutterance \stepcounter{utterance}  
    & & & \multicolumn{2}{p{0.3\linewidth}}{
        \cellcolor[rgb]{0.9,0.9,0.9}{
            \makecell[{{p{\linewidth}}}]{
                \texttt{\tiny{[GM$|$GM]}}
                \texttt{continue} \\
            }
        }
    }
    & & \\ \\

    \theutterance \stepcounter{utterance}  
    & & \multicolumn{4}{p{0.6\linewidth}}{
        \cellcolor[rgb]{0.9,0.9,0.9}{
            \makecell[{{p{\linewidth}}}]{
                \texttt{\tiny{[P1$\langle$GM]}}
                \texttt{Das ist meineBESCHREIBUNG: 5x5{-}Gitter. Diagonale von oben rechts nach unten links und alles darunter (bzw. links unten) ist mit X gefüllt. Zeile 1: 4 leer, 1 X (ganz rechts). Zeile 2: 3 leer, 2 X. Zeile 3: 2 leer, 3 X. Zeile 4: 1 leer, 4 X. Zeile 5: 5 X.} \\
\texttt{FRAGE: Welchen Inhalt haben die Ecken des Rasters (Positionen Zeile 1 Spalte 1, Zeile 1 Spalte 5, Zeile 5 Spalte 1, Zeile 5 Spalte 5)?Beginne deine Antwort mit "ANTWORT:" und füge nichts anderes hinzu.} \\
            }
        }
    }
    & \\ \\

    \theutterance \stepcounter{utterance}  
    & \multicolumn{4}{p{0.6\linewidth}}{
        \cellcolor[rgb]{0.9,0.9,0.9}{
            \makecell[{{p{\linewidth}}}]{
                \texttt{\tiny{[P1$\rangle$GM]}}
                \texttt{ANTWORT: Zeile 1 Spalte 1: leer (▢), Zeile 1 Spalte 5: X, Zeile 5 Spalte 1: X, Zeile 5 Spalte 5: X.} \\
            }
        }
    }
    & & \\ \\

    \theutterance \stepcounter{utterance}  
    & & & \multicolumn{2}{p{0.3\linewidth}}{
        \cellcolor[rgb]{0.9,0.9,0.9}{
            \makecell[{{p{\linewidth}}}]{
                \texttt{\tiny{[GM$|$GM]}}
                \texttt{continue} \\
            }
        }
    }
    & & \\ \\

    \theutterance \stepcounter{utterance}  
    & & \multicolumn{4}{p{0.6\linewidth}}{
        \cellcolor[rgb]{0.9,0.9,0.9}{
            \makecell[{{p{\linewidth}}}]{
                \texttt{\tiny{[P1$\langle$GM]}}
                \texttt{Stelle jetzt eine Frage, um neue Aspekte meines Rasters herauszufinden, die sich möglicherweise von deinem Raster unterscheiden. Beginne mit "FRAGE:" und füge nichts anderes hinzu.} \\
            }
        }
    }
    & \\ \\

    \theutterance \stepcounter{utterance}  
    & \multicolumn{4}{p{0.6\linewidth}}{
        \cellcolor[rgb]{0.9,0.9,0.9}{
            \makecell[{{p{\linewidth}}}]{
                \texttt{\tiny{[P1$\rangle$GM]}}
                \texttt{FRAGE: Wie viele leere Felder (▢) enthält dein Raster insgesamt und wie verteilen sich diese auf die Zeilen?} \\
            }
        }
    }
    & & \\ \\

    \theutterance \stepcounter{utterance}  
    & & & \multicolumn{2}{p{0.3\linewidth}}{
        \cellcolor[rgb]{0.9,0.9,0.9}{
            \makecell[{{p{\linewidth}}}]{
                \texttt{\tiny{[GM$|$GM]}}
                \texttt{continue} \\
            }
        }
    }
    & & \\ \\

    \theutterance \stepcounter{utterance}  
    & & \multicolumn{4}{p{0.6\linewidth}}{
        \cellcolor[rgb]{0.9,0.9,0.9}{
            \makecell[{{p{\linewidth}}}]{
                \texttt{\tiny{[GM$\rangle$P2]}}
                \texttt{ANTWORT: Zeile 1 Spalte 1: leer (▢), Zeile 1 Spalte 5: X, Zeile 5 Spalte 1: X, Zeile 5 Spalte 5: X.} \\
\texttt{FRAGE: Wie viele leere Felder (▢) enthält dein Raster insgesamt und wie verteilen sich diese auf die Zeilen?Beginne deine Antwort mit "ANTWORT:" und füge nichts anderes hinzu.} \\
            }
        }
    }
    & \\ \\

    \theutterance \stepcounter{utterance}  
    & & & \multicolumn{4}{p{0.6\linewidth}}{
        \cellcolor[rgb]{0.9,0.9,0.9}{
            \makecell[{{p{\linewidth}}}]{
                \texttt{\tiny{[GM$\langle$P2]}}
                \texttt{ANTWORT: Mein Raster enthält insgesamt 10 leere Felder. Verteilung auf die Zeilen: Zeile 1: 4, Zeile 2: 3, Zeile 3: 2, Zeile 4: 1, Zeile 5: 0.} \\
            }
        }
    }
     \\ \\

    \theutterance \stepcounter{utterance}  
    & & & \multicolumn{2}{p{0.3\linewidth}}{
        \cellcolor[rgb]{0.9,0.9,0.9}{
            \makecell[{{p{\linewidth}}}]{
                \texttt{\tiny{[GM$|$GM]}}
                \texttt{continue} \\
            }
        }
    }
    & & \\ \\

    \theutterance \stepcounter{utterance}  
    & & \multicolumn{4}{p{0.6\linewidth}}{
        \cellcolor[rgb]{0.9,0.9,0.9}{
            \makecell[{{p{\linewidth}}}]{
                \texttt{\tiny{[GM$\rangle$P2]}}
                \texttt{Stelle jetzt eine Frage, um neue Aspekte meines Rasters herauszufinden, die sich möglicherweise von deinem Raster unterscheiden. Beginne mit "FRAGE:" und füge nichts anderes hinzu.} \\
            }
        }
    }
    & \\ \\

    \theutterance \stepcounter{utterance}  
    & & & \multicolumn{4}{p{0.6\linewidth}}{
        \cellcolor[rgb]{0.9,0.9,0.9}{
            \makecell[{{p{\linewidth}}}]{
                \texttt{\tiny{[GM$\langle$P2]}}
                \texttt{FRAGE: In Zeile 3: Welche Spalten sind genau leer und welche sind X?} \\
            }
        }
    }
     \\ \\

    \theutterance \stepcounter{utterance}  
    & & & \multicolumn{2}{p{0.3\linewidth}}{
        \cellcolor[rgb]{0.9,0.9,0.9}{
            \makecell[{{p{\linewidth}}}]{
                \texttt{\tiny{[GM$|$GM]}}
                \texttt{continue} \\
            }
        }
    }
    & & \\ \\

    \theutterance \stepcounter{utterance}  
    & & \multicolumn{4}{p{0.6\linewidth}}{
        \cellcolor[rgb]{0.9,0.9,0.9}{
            \makecell[{{p{\linewidth}}}]{
                \texttt{\tiny{[P1$\langle$GM]}}
                \texttt{ANTWORT: Mein Raster enthält insgesamt 10 leere Felder. Verteilung auf die Zeilen: Zeile 1: 4, Zeile 2: 3, Zeile 3: 2, Zeile 4: 1, Zeile 5: 0.} \\
\texttt{FRAGE: In Zeile 3: Welche Spalten sind genau leer und welche sind X?Beginne deine Antwort mit "ANTWORT:" und füge nichts anderes hinzu.} \\
            }
        }
    }
    & \\ \\

    \theutterance \stepcounter{utterance}  
    & \multicolumn{4}{p{0.6\linewidth}}{
        \cellcolor[rgb]{0.9,0.9,0.9}{
            \makecell[{{p{\linewidth}}}]{
                \texttt{\tiny{[P1$\rangle$GM]}}
                \texttt{ANTWORT: Zeile 3: Spalte 1: X, Spalte 2: X, Spalte 3: leer (▢), Spalte 4: X, Spalte 5: X.} \\
            }
        }
    }
    & & \\ \\

    \theutterance \stepcounter{utterance}  
    & & & \multicolumn{2}{p{0.3\linewidth}}{
        \cellcolor[rgb]{0.9,0.9,0.9}{
            \makecell[{{p{\linewidth}}}]{
                \texttt{\tiny{[GM$|$GM]}}
                \texttt{continue} \\
            }
        }
    }
    & & \\ \\

    \theutterance \stepcounter{utterance}  
    & & \multicolumn{4}{p{0.6\linewidth}}{
        \cellcolor[rgb]{0.9,0.9,0.9}{
            \makecell[{{p{\linewidth}}}]{
                \texttt{\tiny{[P1$\langle$GM]}}
                \texttt{Stelle jetzt eine Frage, um neue Aspekte meines Rasters herauszufinden, die sich möglicherweise von deinem Raster unterscheiden. Beginne mit "FRAGE:" und füge nichts anderes hinzu.} \\
            }
        }
    }
    & \\ \\

    \theutterance \stepcounter{utterance}  
    & \multicolumn{4}{p{0.6\linewidth}}{
        \cellcolor[rgb]{0.9,0.9,0.9}{
            \makecell[{{p{\linewidth}}}]{
                \texttt{\tiny{[P1$\rangle$GM]}}
                \texttt{FRAGE: Welchen Inhalt haben die Zellen auf der Hauptdiagonalen (Positionen Zeile 1 Spalte 1, Zeile 2 Spalte 2, Zeile 3 Spalte 3, Zeile 4 Spalte 4, Zeile 5 Spalte 5)?} \\
            }
        }
    }
    & & \\ \\

    \theutterance \stepcounter{utterance}  
    & & & \multicolumn{2}{p{0.3\linewidth}}{
        \cellcolor[rgb]{0.9,0.9,0.9}{
            \makecell[{{p{\linewidth}}}]{
                \texttt{\tiny{[GM$|$GM]}}
                \texttt{continue} \\
            }
        }
    }
    & & \\ \\

    \theutterance \stepcounter{utterance}  
    & & \multicolumn{4}{p{0.6\linewidth}}{
        \cellcolor[rgb]{0.9,0.9,0.9}{
            \makecell[{{p{\linewidth}}}]{
                \texttt{\tiny{[GM$\rangle$P2]}}
                \texttt{ANTWORT: Zeile 3: Spalte 1: X, Spalte 2: X, Spalte 3: leer (▢), Spalte 4: X, Spalte 5: X.} \\
\texttt{FRAGE: Welchen Inhalt haben die Zellen auf der Hauptdiagonalen (Positionen Zeile 1 Spalte 1, Zeile 2 Spalte 2, Zeile 3 Spalte 3, Zeile 4 Spalte 4, Zeile 5 Spalte 5)?Beginne deine Antwort mit "ANTWORT:" und füge nichts anderes hinzu.} \\
            }
        }
    }
    & \\ \\

    \theutterance \stepcounter{utterance}  
    & & & \multicolumn{4}{p{0.6\linewidth}}{
        \cellcolor[rgb]{0.9,0.9,0.9}{
            \makecell[{{p{\linewidth}}}]{
                \texttt{\tiny{[GM$\langle$P2]}}
                \texttt{ANTWORT: Hauptdiagonale: Zeile 1 Spalte 1: leer (▢), Zeile 2 Spalte 2: leer (▢), Zeile 3 Spalte 3: X, Zeile 4 Spalte 4: X, Zeile 5 Spalte 5: X.} \\
            }
        }
    }
     \\ \\

    \theutterance \stepcounter{utterance}  
    & & & \multicolumn{2}{p{0.3\linewidth}}{
        \cellcolor[rgb]{0.9,0.9,0.9}{
            \makecell[{{p{\linewidth}}}]{
                \texttt{\tiny{[GM$|$GM]}}
                \texttt{continue} \\
            }
        }
    }
    & & \\ \\

    \theutterance \stepcounter{utterance}  
    & & \multicolumn{4}{p{0.6\linewidth}}{
        \cellcolor[rgb]{0.9,0.9,0.9}{
            \makecell[{{p{\linewidth}}}]{
                \texttt{\tiny{[GM$\rangle$P2]}}
                \texttt{Stelle jetzt eine Frage, um neue Aspekte meines Rasters herauszufinden, die sich möglicherweise von deinem Raster unterscheiden. Beginne mit "FRAGE:" und füge nichts anderes hinzu.} \\
            }
        }
    }
    & \\ \\

    \theutterance \stepcounter{utterance}  
    & & & \multicolumn{4}{p{0.6\linewidth}}{
        \cellcolor[rgb]{0.9,0.9,0.9}{
            \makecell[{{p{\linewidth}}}]{
                \texttt{\tiny{[GM$\langle$P2]}}
                \texttt{FRAGE: Wie viele leere Felder hat deine Zeile 4 und in welchen Spalten befinden sie sich?} \\
            }
        }
    }
     \\ \\

    \theutterance \stepcounter{utterance}  
    & & & \multicolumn{2}{p{0.3\linewidth}}{
        \cellcolor[rgb]{0.9,0.9,0.9}{
            \makecell[{{p{\linewidth}}}]{
                \texttt{\tiny{[GM$|$GM]}}
                \texttt{continue} \\
            }
        }
    }
    & & \\ \\

    \theutterance \stepcounter{utterance}  
    & & \multicolumn{4}{p{0.6\linewidth}}{
        \cellcolor[rgb]{0.9,0.9,0.9}{
            \makecell[{{p{\linewidth}}}]{
                \texttt{\tiny{[P1$\langle$GM]}}
                \texttt{ANTWORT: Hauptdiagonale: Zeile 1 Spalte 1: leer (▢), Zeile 2 Spalte 2: leer (▢), Zeile 3 Spalte 3: X, Zeile 4 Spalte 4: X, Zeile 5 Spalte 5: X.} \\
\texttt{FRAGE: Wie viele leere Felder hat deine Zeile 4 und in welchen Spalten befinden sie sich?Beginne deine Antwort mit "ANTWORT:" und füge nichts anderes hinzu.} \\
            }
        }
    }
    & \\ \\

    \theutterance \stepcounter{utterance}  
    & \multicolumn{4}{p{0.6\linewidth}}{
        \cellcolor[rgb]{0.9,0.9,0.9}{
            \makecell[{{p{\linewidth}}}]{
                \texttt{\tiny{[P1$\rangle$GM]}}
                \texttt{ANTWORT: Zeile 4 hat 1 leeres Feld in Spalte 3.} \\
            }
        }
    }
    & & \\ \\

    \theutterance \stepcounter{utterance}  
    & & & \multicolumn{2}{p{0.3\linewidth}}{
        \cellcolor[rgb]{0.9,0.9,0.9}{
            \makecell[{{p{\linewidth}}}]{
                \texttt{\tiny{[GM$|$GM]}}
                \texttt{continue} \\
            }
        }
    }
    & & \\ \\

    \theutterance \stepcounter{utterance}  
    & & \multicolumn{4}{p{0.6\linewidth}}{
        \cellcolor[rgb]{0.9,0.9,0.9}{
            \makecell[{{p{\linewidth}}}]{
                \texttt{\tiny{[P1$\langle$GM]}}
                \texttt{Stelle jetzt eine Frage, um neue Aspekte meines Rasters herauszufinden, die sich möglicherweise von deinem Raster unterscheiden. Beginne mit "FRAGE:" und füge nichts anderes hinzu.} \\
            }
        }
    }
    & \\ \\

    \theutterance \stepcounter{utterance}  
    & \multicolumn{4}{p{0.6\linewidth}}{
        \cellcolor[rgb]{0.9,0.9,0.9}{
            \makecell[{{p{\linewidth}}}]{
                \texttt{\tiny{[P1$\rangle$GM]}}
                \texttt{FRAGE: Welchen Inhalt hat die Zelle in Zeile 4, Spalte 1?} \\
            }
        }
    }
    & & \\ \\

    \theutterance \stepcounter{utterance}  
    & & & \multicolumn{2}{p{0.3\linewidth}}{
        \cellcolor[rgb]{0.9,0.9,0.9}{
            \makecell[{{p{\linewidth}}}]{
                \texttt{\tiny{[GM$|$GM]}}
                \texttt{continue} \\
            }
        }
    }
    & & \\ \\

    \theutterance \stepcounter{utterance}  
    & & \multicolumn{4}{p{0.6\linewidth}}{
        \cellcolor[rgb]{0.9,0.9,0.9}{
            \makecell[{{p{\linewidth}}}]{
                \texttt{\tiny{[GM$\rangle$P2]}}
                \texttt{ANTWORT: Zeile 4 hat 1 leeres Feld in Spalte 3.} \\
\texttt{FRAGE: Welchen Inhalt hat die Zelle in Zeile 4, Spalte 1?Beginne deine Antwort mit "ANTWORT:" und füge nichts anderes hinzu.} \\
            }
        }
    }
    & \\ \\

    \theutterance \stepcounter{utterance}  
    & & & \multicolumn{4}{p{0.6\linewidth}}{
        \cellcolor[rgb]{0.9,0.9,0.9}{
            \makecell[{{p{\linewidth}}}]{
                \texttt{\tiny{[GM$\langle$P2]}}
                \texttt{ANTWORT: Zeile 4, Spalte 1: X} \\
            }
        }
    }
     \\ \\

    \theutterance \stepcounter{utterance}  
    & & & \multicolumn{2}{p{0.3\linewidth}}{
        \cellcolor[rgb]{0.9,0.9,0.9}{
            \makecell[{{p{\linewidth}}}]{
                \texttt{\tiny{[GM$|$GM]}}
                \texttt{continue} \\
            }
        }
    }
    & & \\ \\

    \theutterance \stepcounter{utterance}  
    & & \multicolumn{4}{p{0.6\linewidth}}{
        \cellcolor[rgb]{0.9,0.9,0.9}{
            \makecell[{{p{\linewidth}}}]{
                \texttt{\tiny{[GM$\rangle$P2]}}
                \texttt{Triff jetzt eine Entscheidung. Was denkst du: Sind dein Raster und das andere beschriebene Raster identisch? Schreibe "ENTSCHEIDUNG: gleiches Raster", wenn du denkst, dass es dasselbe Raster ist, oder "ENTSCHEIDUNG: verschiedene Raster", wenn du denkst, dass es unterschiedliche Raster sind. Füge nichts anderes hinzu.} \\
            }
        }
    }
    & \\ \\

    \theutterance \stepcounter{utterance}  
    & & & \multicolumn{4}{p{0.6\linewidth}}{
        \cellcolor[rgb]{0.9,0.9,0.9}{
            \makecell[{{p{\linewidth}}}]{
                \texttt{\tiny{[GM$\langle$P2]}}
                \texttt{ENTSCHEIDUNG: verschiedene Raster} \\
            }
        }
    }
     \\ \\

    \theutterance \stepcounter{utterance}  
    & & & \multicolumn{2}{p{0.3\linewidth}}{
        \cellcolor[rgb]{0.9,0.9,0.9}{
            \makecell[{{p{\linewidth}}}]{
                \texttt{\tiny{[GM$|$GM]}}
                \texttt{continue} \\
            }
        }
    }
    & & \\ \\

    \theutterance \stepcounter{utterance}  
    & & & \multicolumn{2}{p{0.3\linewidth}}{
        \cellcolor[rgb]{0.9,0.9,0.9}{
            \makecell[{{p{\linewidth}}}]{
                \texttt{\tiny{[GM$|$GM]}}
                \texttt{success} \\
            }
        }
    }
    & & \\ \\

    \theutterance \stepcounter{utterance}  
    & & \multicolumn{4}{p{0.6\linewidth}}{
        \cellcolor[rgb]{0.9,0.9,0.9}{
            \makecell[{{p{\linewidth}}}]{
                \texttt{\tiny{[P1$\langle$GM]}}
                \texttt{ANTWORT: Zeile 4, Spalte 1: X} \\
\texttt{Triff jetzt eine Entscheidung. Was denkst du: Sind dein Raster und das andere beschriebene Raster identisch? Schreibe "ENTSCHEIDUNG: gleiches Raster", wenn du denkst, dass es dasselbe Raster ist, oder "ENTSCHEIDUNG: verschiedene Raster", wenn du denkst, dass es unterschiedliche Raster sind. Füge nichts anderes hinzu.} \\
            }
        }
    }
    & \\ \\

    \theutterance \stepcounter{utterance}  
    & \multicolumn{4}{p{0.6\linewidth}}{
        \cellcolor[rgb]{0.9,0.9,0.9}{
            \makecell[{{p{\linewidth}}}]{
                \texttt{\tiny{[P1$\rangle$GM]}}
                \texttt{ENTSCHEIDUNG: verschiedene Raster} \\
            }
        }
    }
    & & \\ \\

    \theutterance \stepcounter{utterance}  
    & & & \multicolumn{2}{p{0.3\linewidth}}{
        \cellcolor[rgb]{0.9,0.9,0.9}{
            \makecell[{{p{\linewidth}}}]{
                \texttt{\tiny{[GM$|$GM]}}
                \texttt{continue} \\
            }
        }
    }
    & & \\ \\

    \theutterance \stepcounter{utterance}  
    & & & \multicolumn{2}{p{0.3\linewidth}}{
        \cellcolor[rgb]{0.9,0.9,0.9}{
            \makecell[{{p{\linewidth}}}]{
                \texttt{\tiny{[GM$|$GM]}}
                \texttt{success} \\
            }
        }
    }
    & & \\ \\

\end{supertabular}
}

\end{document}
